\documentclass[a4paper, 11pt, titlepage]{article}

\usepackage[utf8]{inputenc}
\usepackage[ngerman]{babel}
\usepackage{graphicx}
\usepackage[colorlinks,linkcolor=black,linktoc=all]{hyperref}

\begin{document}

\begin{titlepage}
\begin{center}

% oberer Teil
\includegraphics[width=0.2\textwidth]{./ico}\\[1cm]

\textsc{\Large Nachhaltigkeit }\\[1.5cm]

% Titel
\rule{\linewidth}{0.5mm} \\[0.4cm]
{ \Huge \bfseries Discamus}\\[0.4cm]
\rule{\linewidth}{0.7mm} \\[1.5cm]

% Autoren
\Large \textbf{Klasse 10/1} \\[1cm]

\large
Manuel \textsc{Bucher}\\
Niko \textsc{Lockenvitz}\\
Edwin \textsc{Brüseke}\\
Julius \textsc{Bittner}


\vfill

% unterer Teil 

{\large \today}

\end{center}
\end{titlepage}
\tableofcontents

\section{Discamus}
\textbf{Discamus} ist ein nachhaltiges Computerprogramm, mithilfe dessen man einfacher die Zuordnungen von Begriffen spielerisch erlernen kann. Es geht dabei in erster Linie um Vokabeln und Ähnliches, beispielsweise die Zuordnungen von Bundesländern oder Staaten zu den entsprechenden Hauptstädten.\\\\
Der Name Discamus rührt vom Inhalt des Programms. Discamus kommt aus dem Lateinischen und bedeutet \glqq Lasst uns lernen\grqq , was sowohl einen positiven, motivierten Grundgedanken als auch einen Gemeinschaftsgedanken beinhaltet.\\\\
Neben den Aspekten der Nachhaltigkeit (s. u.) hat ein Computerprogramm die Vorteile, dass man es nicht betrügen kann, dementsprechend auch sich selbst beim Lernen nicht betrügen kann. Außerdem kann unser Programm eine Auswertung erstellen, ohne dass der Benutzer einen Befehl dazu geben muss. Man sieht als Schüler also auch, wie man sich verbessert hat, welches ungemein zur emotionalen Verbesserung des Schülers beiträgt. Selbstverständlich kann Discamus nicht nur von Schülern, sondern auch von anderen Zielgruppen genutzt werden.\\\\
Die Bedienung erläutern wir im Anhang.

\section{Aspekte der Nachhaltigkeit}
Viele Schüler haben Probleme beim Lernen. Unser Programm jedoch hilft dabei, Schüler beim Lernen zu unterstützen und somit mehr Freizeit zu geben, in denen sie \textendash\ solange der Wille besteht \textendash\ sich nachhaltig engagieren können. Weiterhin ist das Arbeitsmittel Computer ein weiterer Ansporn für junge Menschen, zu lernen, da Computer (im Gegensatz zu Büchern oder Tageszeitungen) als nicht mehr modern gelten und nicht so gerne benutzt werden.\\\\
Die Schüler haben, wenn sie ordentlich gelernt haben, im Endeffekt bessere Noten und dadurch bessere Perspektiven für ihre eigene Zukunft. Wer als erwachsener Mensch im Leben besser da steht, kann beispielsweise durch eine gute finanzielle Lage Spenden abgeben, um Menschen zu helfen, die das Geld dringend benötigen, es aber nicht selber aufwenden können.\\\\
Im Gegensatz zu herkömmlichen Lernmethoden bietet Discamus weiterhin den Vorteil, dass ein Rechner im Gegensatz zum Menschen kein Nachsehen hat und der Anwender keine Möglichkeit zum Schummeln hat. Außerdem kann unser Lernprogramm die Eingaben auswerten, langfristig die Veränderungen speichern und dem Anwender anzeigen, sodass eine Verbesserung ein weiterer Motivationsfaktor für Schüler sein kann. Nicht zuletzt nimmt das Lernen außerdem weniger Zeit in Anspruch und schützt die Umwelt, da weder Papier noch Tinte verbraucht oder für überflüssige Notizen verschwendet werden müssen.\\\\
Es gibt zwar auch noch andere Computerlernprogramme. Allerdings sind diese fast immer auf ein bestimmtes Gebiet spezialisiert und haben vordefinierte Aufgaben, sodass sie nur für spezielle Gebiete einsetzbar sind. Discamus ist ebenso effektiv wie andere Lernprogramme, aber simpler zu bedienen und einfacher zu verstehen als die meisten Konkurrenten.\\\\
Unser Programm verbraucht auch viel weniger Speicherplatz als die meisten Lernprogramme. Es kann dank des Quellcodes, der in der Programmiersprache Python geschrieben ist, auf allen gängigen Betriebssystemen ohne Probleme ausgeführt werden. Aufgrund dessen, dass der Quellcode frei einsehbar ist, ist auch Sicherheit garantiert. Weiterhin ist unser Programm kostenlos, frei von Werbung und bietet Datenschutz. Wer einen Ansprechpartner sucht, kann uns (im Rahmen der Schule) direkt ansprechen, ohne den langwierigen Weg über E-Mail oder Web-Formulare wählen zu müssen. Auch die Verbreitung über USB-Sticks ist einfach möglich. Unser Programm muss nicht einmal installiert werden, sondern kann direkt ausgeführt werden.

\section{Fazit}

Man kann sagen, dass das Programm dem lernenden Schüler hilft. Wenn der Schüler Interesse zeigt, wird er mit Hilfe des Programms sogar besser. Neben diesen Aspekten bietet es aber auch im Vergleich zu ähnlichen Produkten Vorteile. Gerade in der heutigen Zeit, wo Datenschutz groß geschrieben wird, stellt unser Programm einen riesigen Vorteil dar. Auch in der Zukunft wird unser Programm durch die einfache Distribution, die starke Simplizität und insbesondere den groß geschriebenen Datenschutz bestehen.


\section{Arbeitseinteilung}

Für die Entwicklung unseres Programms haben wir die Arbeit wie folgt eingeteilt: Während Edwin die grafische Benutzeroberfläche eingerichtet und die Betriebsanweisung (Anhang) zum Programm erarbeitet hat, hat Niko die einzelnen Programmteile, die Manuel auf Nikos Beschreibung hin programmiert hat, sinnvoll zusammengestellt. Julius hat diesen Bericht zur Projektierung fertiggestellt. Letztendlich hat aber jeder an jedem Teil der Gestaltung des Programms mitgewirkt.\\\\
Zur einfacheren Bewertung zeigen wir in der folgenden Tabelle unsere individuellen Arbeitszeiten am Projekt auf, die Zeiten sind in Stunden angegeben. Der hohe Unterschied in der Gesamtzahl der Stunden kommt daher zustande, dass einige Arbeiten sich als schwieriger und langwieriger herausstellten als ursprünglich vermutet, was vorher nicht abzuschätzen war. Es sind alle Arbeitstage angegeben; an fehlenden Tagen haben wir nicht aktiv zur Fertigstellung des Projekts beigetragen.
\begin{table}[htbp]
\centering
\begin{tabular}{|c||c|c|c|c|}
\hline
Datum & Niko & Edwin & Manuel & Julius \\
\hline
\hline
05.12.2014 & 2 & 2 & 2 & 2 \\
\hline
06.12.2014 & 2 & 0 & 1 & 0 \\
\hline
07.12.2014 & 0 & 0 & 1 & 0 \\
\hline
08.12.2014 & 0 & 2 & 3 & 0 \\
\hline
09.12.2014 & 0 & 1 & 0{,}5 & 3 \\
\hline
10.12.2014 & 0{,}5 & 0 & 0 & 0{,}5 \\
\hline
11.12.2014 & 2 & 1 & 2 & 0 \\
\hline
12.12.2014 & 1 & 1 & 0 & 0 \\
\hline
13.12.2014 & 1 & 0 & 0 & 0 \\
\hline
14.12.2014 & 2 & 0 & 2 & 0 \\
\hline
15.12.2014 & 0 & 0 & 2 & 0 \\
\hline
16.12.2014 & 1 & 0 & 0{,}5 & 0{,}5 \\
\hline
17.12.2014 & 0 & 0 & 1{,}5 & 0{,}5 \\
\hline
19.12.2014 & 0 & 0 & 2 & 0 \\
\hline
20.12.2014 & 0 & 0 & 1 & 0 \\
\hline
21.12.2014 & 2 & 0 & 2 & 0 \\
\hline
30.12.2014 & 0 & 0 & 1 & 0 \\
\hline
01.01.2015 & 0 & 0 & 1 & 0 \\
\hline
02.01.2015 & 1 & 0 & 0 & 0 \\
\hline
03.01.2015 & 4 & 0 & 0 & 0 \\
\hline
04.01.2015 & 2{,}5 & 0 & 2 & 0 \\
\hline
06.01.2015 & 1 & 0 & 0 & 0 \\
\hline
09.01.2015 & 0 & 0 & 1 & 0 \\
\hline
10.01.2015 & 1{,}5 & 0 & 1 & 0 \\
\hline
11.01.2015 & 0{,}5 & 0 & 0 & 0 \\
\hline
17.01.2015 & 0 & 0 & 1 & 0 \\
\hline
19.01.2015 & 0 & 0 & 1 & 0 \\
\hline
20.01.2015 & 0 & 0 & 1 & 0 \\
\hline
23.02.2015 & 1 & 0 & 1 & 0 \\
\hline
24.02.2015 & 3 & 3 & 3 & 3 \\
\hline
22.03.2015 & 1 & 0 & 0 & 3 \\
\hline
\hline
gesamt & 29 & 10 & 33{,}5 & 13 \\
\hline
\end{tabular}
\caption{Arbeitszeiten}
\end{table}

\end{document}
